% from ega4-17-fr.tex
\oldpage[IV-4]{62}

\begin{enumerate}
  \item[(v)] Soient $f:X\to Y$, $g:Y\to Z$ deux morphismes; si $g$ est localement de type fini et si $g\circ f$ est non ramifié, alors $f$ est non ramifié.
\end{enumerate}
\end{proposition}

Cela résulte aussitôt de (I, 4.3) et de (17.1.3).

\begin{proposition}[17.3.4]
\label{IV.17.3.4}
Soient $f:X\to Y$, $g:Y\to Z$ deux morphismes, et supposons $g$ non ramifié. Alors, si $g\circ f$ est lisse (resp. non ramifié, resp. étale), il en est de même de $f$.
\end{proposition}

En effet, comme $g$ et $g\circ f$ sont localement de présentation finie, il en est de même de $f$ (I, 4.4.3, (v)); la conclusion résulte donc de (17.1.4) et (17.1.3, (v)).

\begin{corollary}[17.3.5]
\label{IV.17.3.5}
Supposons $g$ étale; alors, pour que $f$ soit lisse (resp. non ramifié, resp. étale), il faut et il suffit que $g\circ f$ le soit.
\end{corollary}

Cela résulte de (17.3.4) et de (17.3.3, (iii)).

\begin{proposition}[17.3.6]
\label{IV.17.3.6}
Soient $g:Y\to S$, $h:X\to S$ deux morphismes localement de présentation finie. Pour qu'un $S$-morphisme $f:X\to Y$ soit non ramifié, il faut et il suffit que l'homomorphisme canonique (16.4.19)
\[
  f^*(\Omega_{Y/S}^1)\longrightarrow\Omega_{X/S}^1
\]
soit surjectif.
\end{proposition}

Comme $f$ est alors localement de présentation finie (I, 4.4.3, (v)), la proposition résulte aussitôt de (17.2.2).

\begin{definition}[17.3.7]
\EGAReaderTarget{ega4.c04.def.17-3-7}{IV.17.3.7}
Soit $f:X\to Y$ un morphisme. On dit que $f$ est \emph{lisse} (resp. \emph{non ramifié}, resp. \emph{étale}) \emph{en un point} $x\in X$ s'il existe un voisinage ouvert $U$ de $x$ dans $X$ tel que la restriction $f|U$ soit un morphisme lisse (resp. non ramifié, resp. étale) de $U$ dans $Y$.
\end{definition}

On dit encore alors que $X$ est lisse (resp. non ramifié, resp. étale) sur $Y$ au point $x$.

Compte tenu de la remarque (17.3.2, (iii)), il revient au même de dire que $f$ est un morphisme lisse (resp. non ramifié, resp. étale) ou qu'il est lisse (resp. non ramifié, resp. étale) en chaque point de $X$.

Il est clair que l'ensemble des points de $X$ où un morphisme $f:X\to Y$ est lisse (resp. non ramifié, resp. étale) est \emph{ouvert} dans $X$.

\begin{proposition}[17.3.8]
\label{IV.17.3.8}
Pour tout préschéma $Y$ et tout $\sh{O}_Y$-Module localement libre $\sh{E}$ de type fini, le préschéma fibré vectoriel $\mathbf{V}(\sh{E})$ (II, 1.7.8) est un $Y$-préschéma lisse.
\end{proposition}

En effet (17.3.2, (iii)), on peut se borner au cas où $Y=\Spec(A)$ est affine et $\mathbf{V}(\sh{E})=\Spec(A[T_1,\ldots,T_r])$; comme $A[T_1,\ldots,T_r]$ est une $A$-algèbre formellement lisse pour les topologies discrètes (0, 19.3.2) et de présentation finie, cela démontre la proposition (17.3.2, (ii)).

\begin{corollary}[17.3.9]
\label{IV.17.3.9}
Sous les hypothèses de (17.3.8), le préschéma fibré projectif $\mathbf{P}(\sh{E})$ (II, 4.1.1) est un $Y$-préschéma lisse.
\end{corollary}

On peut encore se borner au cas où $Y=\Spec(A)$ est affine et $\mathbf{P}(\sh{E})=\mathbf{P}^r_Y$. On sait alors (II, 2.3.14) qu'on a un recouvrement ouvert fini de $\mathbf{P}^r_A$ en prenant les $D_+(T_i)$ ($0\leq i\leq r$), égaux respectivement aux spectres des anneaux $S_{(f)}$, où l'on remplace $S$ par $A[T_0,T_1,\ldots,T_r]$ et $f$ par $T_i$; mais il résulte aussitôt de la définition de $S_{(f)}$ (II, 2.2.1) que cet anneau, dans le cas considéré, est isomorphe à $A[T_0,\ldots,T_{i-1},T_{i+1},\ldots,T_r]$; donc le corollaire résulte de (17.3.8).

